% Compile from the Pogotrack repository root with LuaLaTeX or pdfLaTeX.
% This file is a blank template for documenting Pogobot experiments.

\documentclass[11pt,a4paper]{article}

\usepackage[margin=22mm,headheight=28pt,footskip=16mm]{geometry}
\usepackage[T1]{fontenc}
\usepackage[utf8]{inputenc}
\usepackage[english]{babel}
\usepackage{lmodern}
\usepackage{microtype}
\usepackage[table]{xcolor}
\usepackage{array}
\usepackage{graphicx}
\usepackage{booktabs}
\usepackage{enumitem}
\usepackage{tabularx}
\usepackage{caption}
\usepackage{float}
\usepackage{listings}
\usepackage{tcolorbox}
\usepackage{tikz}
\usepackage{fancyhdr}
\usepackage[hidelinks]{hyperref}
\usepackage{url}

\tcbuselibrary{breakable,listings}

\graphicspath{{img/}}
\setlength{\parindent}{0pt}
\setlength{\parskip}{0.55em}
\renewcommand{\familydefault}{\sfdefault}
\setcounter{tocdepth}{1}

% ---------- Colour palette ----------
\definecolor{PogoGreen}{HTML}{126B5E}
\definecolor{PogoDark}{HTML}{12332E}
\definecolor{PogoLight}{HTML}{EAF4F1}
\definecolor{PogoOrange}{HTML}{E47735}
\definecolor{PogoRed}{HTML}{B53B37}
\definecolor{PogoGray}{HTML}{4D5654}
\definecolor{CodeBackground}{HTML}{F4F7F6}

\hypersetup{
  pdftitle={Tutorial: Pogobot experiments},
  colorlinks=true,
  linkcolor=PogoGreen,
  urlcolor=PogoGreen,
  citecolor=PogoGreen
}

\captionsetup{
  font=small,
  labelfont={bf,color=PogoGreen},
  textfont=it,
  skip=7pt
}

\setlist[itemize]{leftmargin=1.45em,itemsep=0.35em,topsep=0.35em}
\setlist[enumerate]{leftmargin=1.65em,itemsep=0.55em,topsep=0.35em,label=\textcolor{PogoGreen}{\arabic*.}}

\pagestyle{fancy}
\fancyhf{}
\lhead{\textcolor{PogoGreen}{\textbf{Tutorial: Pogobot experiments}}}
\rhead{\textcolor{PogoGray}{September 2026}}
\cfoot{\textcolor{PogoGray}{\thepage}}
\renewcommand{\headrulewidth}{0.5pt}
\renewcommand{\headrule}{\hbox to\headwidth{\color{PogoGreen}\leaders\hrule height \headrulewidth\hfill}}

\fancypagestyle{titlepage}{
  \fancyhf{}
  \cfoot{\textcolor{PogoGray}{\thepage}}
  \renewcommand{\headrulewidth}{0pt}
}

\newcommand{\important}[1]{%
  \begin{tcolorbox}[
    colback=PogoLight,
    colframe=PogoGreen,
    boxrule=0.7pt,
    arc=2mm,
    left=4mm,right=4mm,top=2.5mm,bottom=2.5mm
  ]
  \textcolor{PogoDark}{\textbf{Important.}}~#1
  \end{tcolorbox}%
}

\newcommand{\warning}[1]{%
  \begin{tcolorbox}[
    colback=PogoOrange!10,
    colframe=PogoOrange,
    boxrule=0.8pt,
    arc=2mm,
    left=4mm,right=4mm,top=2.5mm,bottom=2.5mm
  ]
  \textcolor{PogoOrange!85!black}{\textbf{Attention.}}~#1
  \end{tcolorbox}%
}

\newtcblisting{codeblock}{
  listing only,
  breakable,
  colback=CodeBackground,
  colframe=PogoGreen!55,
  boxrule=0.5pt,
  arc=1.5mm,
  left=3mm,right=3mm,top=2mm,bottom=2mm,
  listing options={
    basicstyle=\ttfamily\footnotesize,
    columns=fullflexible,
    keepspaces=true,
    breaklines=true,
    showstringspaces=false,
    tabsize=2
  }
}

\newcommand{\stepimage}[3]{%
  \begin{figure}[H]
    \centering
    \includegraphics[width=#1\linewidth]{#2}
    \caption{#3}
  \end{figure}%
}

\begin{document}

% ---------- Title block ----------
\begin{tikzpicture}[remember picture,overlay]
  \fill[PogoGreen] (current page.north west) rectangle ([yshift=-42mm]current page.north east);
  \fill[PogoOrange] ([yshift=-42mm]current page.north west) rectangle ([yshift=-44mm]current page.north east);
\end{tikzpicture}
\vspace*{6mm}
{\color{white}\fontsize{23}{29}\selectfont\textbf{Tutorial: Pogobot experiments}\par}
\vspace{25mm}

\thispagestyle{titlepage}

\tableofcontents
\vspace{0.5em}
\hrule height 0.7pt
\vspace{0.8em}

% ---------- Template structure ----------

\section{Introduction}
\subsection{Purpose and scope}
% Describe the purpose of the experiment and the intended reader.

\subsection{Experiment overview}
% Introduce the scientific question, setup, and main stages of the analysis.

\section{Installation}
\subsection{System requirements}
% Record Python, operating-system, and external-tool requirements.

\subsection{Environment setup}
% Add installation commands and instructions for activating the environment.

\section{Repository and experiment structure}
% Describe the relevant repository directories and the expected input files.

\subsection{Input data layout}
% Add the folder tree and explain the naming conventions.

\section{Core analysis: one experiment}
\subsection{Standard tracking}
% Explain how to run the dynamics-video analysis and describe its output.

\subsection{Parallel processing of one video}
% Explain workers, warm-up frames, chunk boundaries, and validation.

\subsection{RGB-ID tracking}
% Explain the RGB-ID video, configuration, output, and limitations.

\subsection{Decoding and merging controller data}
% Explain the merge inputs, timing offset, and final dataset.

\subsection{Generating summary figures}
% Explain how to create and interpret PNG/PDF summary figures.

\section{Configuration}
% Document the YAML parameters that are specific to this experiment.

\section{Config-driven workflow}
\subsection{Configuring an experiment}
% Add a pipeline.yaml example.

\subsection{Running all stages}
% Explain the complete command and the order of operations.

\subsection{Running one stage or one experiment}
% Add commands for batch, RGB-ID, merge, slides, and dry-run modes.

\subsection{Multiprocessing strategy}
% Explain per-video workers versus concurrent experiment workers.

\section{Outputs and interpretation}
% Describe generated files, columns, units, and physical meanings.

\section{Troubleshooting}
% Add common errors, diagnostic commands, and validation checks.

\section{Final remarks}
% Add acknowledgements, contact information, and references.

\end{document}
